\section{Conclusion}\label{sec:conclusion}

This work extends the IJCNN precursor on technical indicators, Information Gain feature selection, and GA-based hyperparameter optimization for Mini-Index futures prediction. The original practical pipeline is retained: a compact set of technical descriptors is constructed from intraday WIN data, informative attributes are selected, and nonlinear models are optimized before final testing. The present study expands that design through a chronological train/validation/test split, four model families, five budget-matched optimizers, two explicit validation objectives, and a locked economic evaluation.

The results show that a single validation score cannot define a universal model ranking. RF produces the strongest mean locked-test accuracy, MCC, and $F_1$ under the compound MCC/$F_1$ fitness, whereas MLP leads the predictive results under weighted accuracy fitness. The economic analysis reaches a related but distinct conclusion: MLP has the largest mean profit under weighted accuracy fitness, while RF has the smallest mean drawdown. These findings support the use of controlled metaheuristic optimization, but they also show why the selection objective, reported predictive metric, and economic endpoint must be stated separately.

The direct comparison between fitness functions remains descriptive because only three seeds are common to every model--optimizer cell. The statistical analysis is correspondingly limited to exploratory within-protocol comparisons. Future work should complete a preregistered common 20-seed design, retain the same temporal split and evaluation budget, and repeat the analysis across additional liquid futures markets and market regimes. Such an extension would test the stability of the present results while preserving the transparent workflow established here.
